\documentclass[11pt]{article}
\usepackage{amsmath, amssymb, geometry, parskip}
\geometry{letterpaper, margin=1in}

\title{\textbf{Homework 2.5: Differential Equations - Exact Equations \& Integrating Factors}}
\author{Tashfeen Omran}
\date{September 2026}

\begin{document}

\maketitle

\section*{Part 1: Comprehensive Introduction, Context, and Prerequisites}

\subsection*{Core Concepts}
The topic at hand is the theory and solution methodology for First-Order Ordinary Differential Equations (ODEs), specifically focusing on \textbf{Exact Equations} and the use of \textbf{Integrating Factors}.
A first-order ODE can frequently be rewritten in the differential form:
\[ M(x, y) \, dx + N(x, y) \, dy = 0 \]
where $M(x, y)$ and $N(x, y)$ are functions of two variables.

An equation in this form is called \textbf{exact} if there exists a continuously differentiable function $F(x, y)$ such that the total differential of $F(x, y)$ matches the left-hand side of the ODE. By the multivariable chain rule, the total differential of $F$ is:
\[ dF = \frac{\partial F}{\partial x} \, dx + \frac{\partial F}{\partial y} \, dy \]
If $dF = M \, dx + N \, dy$, then the ODE becomes $dF = 0$. Integrating both sides yields the implicit solution:
\[ F(x, y) = C \]
where $C$ is an arbitrary constant. For this to happen, the cross-partial derivatives must be equal: $\frac{\partial M}{\partial y} = \frac{\partial N}{\partial x}$. 

When $\frac{\partial M}{\partial y} \neq \frac{\partial N}{\partial x}$, the equation is not exact. However, we can often multiply the entire equation by a special function called an \textbf{integrating factor}, denoted as $\mu(x)$ or $\mu(y)$ or $\mu(x, y)$, which forces the equation to become exact.

\subsection*{Intuition and Derivation}
The intuition stems directly from Multivariable Calculus and Clairaut's Theorem on the equality of mixed partial derivatives. If $M = \frac{\partial F}{\partial x}$ and $N = \frac{\partial F}{\partial y}$, then differentiating $M$ with respect to $y$ gives $\frac{\partial^2 F}{\partial y \partial x}$, and differentiating $N$ with respect to $x$ gives $\frac{\partial^2 F}{\partial x \partial y}$. Clairaut's Theorem guarantees that for smooth functions, $\frac{\partial^2 F}{\partial y \partial x} = \frac{\partial^2 F}{\partial x \partial y}$. Thus, the test for exactness is simply checking if $M_y = N_x$.

If it is not exact, we seek a $\mu$ such that:
\[ \mu M \, dx + \mu N \, dy = 0 \]
is exact. This requires:
\[ \frac{\partial}{\partial y} (\mu M) = \frac{\partial}{\partial x}(\mu N) \]
Applying the product rule yields a partial differential equation for $\mu$. By assuming $\mu$ is a function of only $x$ (or only $y$), this PDE simplifies into a separable first-order ODE for $\mu$, which leads to the exponential formulas used to find integrating factors.

\subsection*{Historical Context and Motivation}
The concept of exact differentials and integrating factors was developed in the 18th century by mathematical titans such as Leonhard Euler and Alexis Clairaut. Mathematicians were attempting to describe physical systems---like planetary motion and fluid dynamics---where equations naturally arose in the form $M \, dx + N \, dy = 0$. Euler formalized the notion of an integrating factor in 1739, realizing that even when an equation was not immediately integrable (not an exact differential), a well-chosen multiplier could make it so, drastically expanding the families of solvable physical problems.

\subsection*{Key Formulas}
\begin{itemize}
    \item \textbf{Standard Form:} $M(x, y) \, dx + N(x, y) \, dy = 0$
    \item \textbf{Test for Exactness:} $M_y = N_x \quad \text{where} \quad M_y = \frac{\partial M}{\partial y}, \; N_x = \frac{\partial N}{\partial x}$
    \item \textbf{Integrating Factor} $\mu(x)$: If $\frac{M_y - N_x}{N}$ is a function of $x$ alone, then:
    \[ \mu(x) = e^{\int \frac{M_y - N_x}{N} \, dx} \]
    \item \textbf{Integrating Factor} $\mu(y)$: If $\frac{N_x - M_y}{M}$ is a function of $y$ alone, then:
    \[ \mu(y) = e^{\int \frac{N_x - M_y}{M} \, dy} \]
\end{itemize}

\subsection*{Prerequisites}
To master this topic, you must be proficient in:
\begin{enumerate}
    \item \textbf{Partial Differentiation:} Treating one variable as a constant while differentiating with respect to the other.
    \item \textbf{Integration Techniques:} Fundamental integration, substitution, and integration by parts.
    \item \textbf{Exponential and Logarithmic Properties:} Specifically $e^{\ln(f(x))} = f(x)$ and $e^{a \ln(x)} = x^a$.
\end{enumerate}

\newpage
\section*{Part 2: Detailed Homework Solutions}

\subsection*{Problem 1}
\textbf{Identify the equation as separable, linear, exact, or having an integrating factor that is a function of either $x$ or $y$ alone.}
\[ (4x + x^{-1}y) \, dx + (xy^3 - 1) \, dy = 0 \]

\textbf{Step 1: Define $M$ and $N$ and test for exactness.}
\begin{align*}
    M(x, y) &= 4x + x^{-1}y \implies M_y = \frac{\partial}{\partial y}(4x + x^{-1}y) = x^{-1} \\
    N(x, y) &= xy^3 - 1 \implies N_x = \frac{\partial}{\partial x}(xy^3 - 1) = y^3
\end{align*}
Since $M_y \neq N_x$ ($x^{-1} \neq y^3$), the equation is \textbf{not exact}.

\textbf{Step 2: Check for an integrating factor $\mu(x)$.} Calculate the expression $\frac{M_y - N_x}{N}$:
\[ \frac{M_y - N_x}{N} = \frac{x^{-1} - y^3}{xy^3 - 1} \]
To simplify, factor out $-x^{-1}$ from the numerator:
\[ \frac{M_y - N_x}{N} = \frac{-x^{-1}(-1 + xy^3)}{xy^3 - 1} = \frac{-x^{-1}(xy^3 - 1)}{xy^3 - 1} = -x^{-1} = -\frac{1}{x} \]
Since this result depends \textit{only} on $x$, there exists an integrating factor $\mu(x)$.

\textbf{Step 3: Calculate the integrating factor.}
\[ \mu(x) = e^{\int -\frac{1}{x} \, dx} = e^{-\ln|x|} = x^{-1} \]
This is an integrating factor not equal to a constant.

\textbf{Final Answer:} The correct option is \textbf{B. has an integrating factor $\mu(x)$ or $\mu(y)$ not equal to a constant.}

\hrulefill

\subsection*{Problem 2}
\textbf{Identify the equation as separable, linear, exact, or having an integrating factor that is a function of either $x$ or $y$ alone.}
\[ (-4x^3 + 3y) \, dx + (3x + 2y) \, dy = 0 \]

\textbf{Step 1: Define $M$ and $N$ and test for exactness.}
\begin{align*}
    M(x, y) &= -4x^3 + 3y \implies M_y = 3 \\
    N(x, y) &= 3x + 2y \implies N_x = 3
\end{align*}
Since $M_y = N_x = 3$, the equation is exact.

\textbf{Final Answer:} The correct option is \textbf{D. exact.}

\hrulefill

\subsection*{Problem 3}
\textbf{Solve the equation:} $(2xy) \, dx + (2y - 3x^2) \, dy = 0$

\textbf{Step 1: Test for exactness.}
\begin{align*}
    M &= 2xy \implies M_y = 2x \\
    N &= 2y - 3x^2 \implies N_x = -6x
\end{align*}
$M_y \neq N_x$, so it is not exact.

\textbf{Step 2: Find an integrating factor.} Let's check $\frac{M_y - N_x}{N}$:
\[ \frac{2x - (-6x)}{2y - 3x^2} = \frac{8x}{2y - 3x^2} \quad \text{(Not a function of $x$ alone)} \]
Let's check $\frac{N_x - M_y}{M}$:
\[ \frac{-6x - 2x}{2xy} = \frac{-8x}{2xy} = -\frac{4}{y} \]
This is a function of $y$ alone, so $\mu(y)$ exists.
\[ \mu(y) = e^{\int -\frac{4}{y} \, dy} = e^{-4\ln|y|} = y^{-4} \]
\textbf{What is an integrating factor for the equation?} $\mu(y) = y^{-4}$.

\textbf{Step 3: Multiply ODE by $\mu(y)$ and solve.}
\[ (2xy^{-3}) \, dx + (2y^{-3} - 3x^2y^{-4}) \, dy = 0 \]
Now $M' = 2xy^{-3}$ and $N' = 2y^{-3} - 3x^2y^{-4}$. Integrate $M'$ with respect to $x$:
\[ F(x, y) = \int 2xy^{-3} \, dx = x^2y^{-3} + g(y) \]
Differentiate $F$ with respect to $y$ and set equal to $N'$:
\[ \frac{\partial F}{\partial y} = -3x^2y^{-4} + g'(y) = 2y^{-3} - 3x^2y^{-4} \]
\[ g'(y) = 2y^{-3} \implies g(y) = \int 2y^{-3} \, dy = -y^{-2} \]
So, $F(x, y) = x^2y^{-3} - y^{-2}$. The implicit solution is $x^2y^{-3} - y^{-2} = C$.

\textbf{Step 4: Check for lost or gained solutions.} We multiplied the original equation by $y^{-4}$, which assumes $y \neq 0$. Let's check if $y = 0$ was a solution to the original equation:
\[ (2x(0)) \, dx + (2(0) - 3x^2)(0) = 0 \implies 0 = 0 \]
Since $y = 0$ satisfies the original equation but cannot be plugged into the general solution (due to division by zero), the solution $y = 0$ was lost.

\textbf{Final Answers:}
\begin{itemize}
    \item Integrating factor: $\mu(y) = y^{-4}$
    \item Implicit solution: $x^2y^{-3} - y^{-2} = C$
    \item Lost/Gained: \textbf{the solution $y = 0$ was lost}
\end{itemize}

\hrulefill

\subsection*{Problem 4}
\textbf{Solve the equation:} $(3x^3 + xy) \, dx + (x^3y - x^2) \, dy = 0$

\textbf{Step 1: Test for exactness.}
\begin{align*}
    M &= 3x^3 + xy \implies M_y = x \\
    N &= x^3y - x^2 \implies N_x = 3x^2y - 2x
\end{align*}
$M_y \neq N_x$, not exact.

\textbf{Step 2: Find integrating factor.}
\[ \frac{M_y - N_x}{N} = \frac{x - (3x^2y - 2x)}{x^2(xy - 1)} = \frac{3x - 3x^2y}{x^2(xy - 1)} = \frac{-3x(xy - 1)}{x^2(xy - 1)} = -\frac{3}{x} \]
Since this depends only on $x$, $\mu(x) = e^{\int -3/x \, dx} = x^{-3}$.

\textbf{Step 3: Multiply ODE by $\mu(x)$ and solve.}
\[ (3 + x^{-2}y) \, dx + (y - x^{-1}) \, dy = 0 \]
Integrate $M'$ with respect to $x$:
\[ F(x, y) = \int (3 + x^{-2}y) \, dx = 3x - x^{-1}y + g(y) \]
Differentiate with respect to $y$ and equate to $N'$:
\[ \frac{\partial F}{\partial y} = -x^{-1} + g'(y) = y - x^{-1} \]
\[ g'(y) = y \implies g(y) = \frac{1}{2}y^2 \]
So, $F(x, y) = 3x - x^{-1}y + \frac{1}{2}y^2$. The solution is $3x - x^{-1}y + \frac{1}{2}y^2 = C$.

\textbf{Step 4: Check for lost or gained solutions.} We multiplied by $x^{-3}$, assuming $x \neq 0$. If $x = 0$ in the original ODE:
\[ (0 + 0) \, dx + (0 - 0) \, dy = 0 \implies 0 = 0 \]
The trivial solution $x = 0$ satisfies the original ODE but is undefined in our implicit solution due to the $x^{-1}$ term. Therefore, $x = 0$ was lost.

\textbf{Final Answers:}
\begin{itemize}
    \item Implicit solution: $3x - x^{-1}y + \frac{1}{2}y^2 = C$
    \item Lost/Gained: \textbf{the solution $x = 0$ was lost}
\end{itemize}

\hrulefill

\subsection*{Problem 5}
\textbf{Solve the equation:} $(3 + 2x^{-1}y^2 + 2x^{-1}y^3) \, dx + (2y + 3y^2) \, dy = 0$

\textbf{Step 1: Test for exactness.}
\begin{align*}
    M_y &= 4x^{-1}y + 6x^{-1}y^2 \\
    N_x &= 0
\end{align*}
$M_y \neq N_x$, not exact.

\textbf{Step 2: Find integrating factor.}
\[ \frac{M_y - N_x}{N} = \frac{4x^{-1}y + 6x^{-1}y^2}{2y + 3y^2} = \frac{2x^{-1}y(2 + 3y)}{y(2 + 3y)} = 2x^{-1} = \frac{2}{x} \]
\[ \mu(x) = e^{\int \frac{2}{x} \, dx} = x^2 \]

\textbf{Step 3: Multiply ODE by $\mu(x)$ and solve.}
\[ (3x^2 + 2xy^2 + 2xy^3) \, dx + (2x^2y + 3x^2y^2) \, dy = 0 \]
Here, it's easier to integrate $N'$ with respect to $y$:
\[ F(x, y) = \int (2x^2y + 3x^2y^2) \, dy = x^2y^2 + x^2y^3 + h(x) \]
Differentiate with respect to $x$ and equate to $M'$:
\[ \frac{\partial F}{\partial x} = 2xy^2 + 2xy^3 + h'(x) = 3x^2 + 2xy^2 + 2xy^3 \]
\[ h'(x) = 3x^2 \implies h(x) = x^3 \]
So, $F(x, y) = x^3 + x^2y^2 + x^2y^3$. Solution: $x^3 + x^2y^2 + x^2y^3 = C$.

\textbf{Step 4: Check for lost or gained solutions.} We multiplied by $x^2$. This means the new equation allows $x = 0$. If we plug $x = 0$ into the new ODE: $0 \, dx + 0 \, dy = 0$, so it is a solution. However, in the original ODE, the terms $2x^{-1}y^2$ are undefined at $x = 0$. Therefore, $x = 0$ was not a valid solution originally, but it became one after multiplication. Thus, $x = 0$ was gained.

\textbf{Final Answers:}
\begin{itemize}
    \item Implicit solution: $x^3 + x^2y^2 + x^2y^3 = C$
    \item Lost/Gained: \textbf{the solution $x = 0$ was gained}
\end{itemize}

\hrulefill

\subsection*{Problem 6}
\textbf{Solve the equation:} $(2xy^4 + 1) \, dx + (4x^2y^3 - 2y^{-1}) \, dy = 0$

\textbf{Step 1: Test for exactness.}
\begin{align*}
    M_y &= \frac{\partial}{\partial y}(2xy^4 + 1) = 8xy^3 \\
    N_x &= \frac{\partial}{\partial x}(4x^2y^3 - 2y^{-1}) = 8xy^3
\end{align*}
$M_y = N_x$. The equation is already exact. $\mu = 1$.

\textbf{Step 2: Solve the exact equation.} Integrate $M$ with respect to $x$:
\[ F(x, y) = \int (2xy^4 + 1) \, dx = x^2y^4 + x + g(y) \]
Differentiate with respect to $y$:
\[ \frac{\partial F}{\partial y} = 4x^2y^3 + g'(y) = 4x^2y^3 - 2y^{-1} \]
\[ g'(y) = -2y^{-1} \implies g(y) = -2\ln|y| \]
Solution: $x^2y^4 + x - 2\ln|y| = C$.

\textbf{Step 3: Check for lost or gained solutions.} Because we only multiplied by the trivial integrating factor $\mu = 1$, we neither assumed any variables were non-zero nor introduced any new roots.

\textbf{Final Answers:}
\begin{itemize}
    \item Implicit solution: $x^2y^4 + x - 2\ln|y| = C$
    \item Lost/Gained: \textbf{no solutions were lost (nor gained)}
\end{itemize}

\hrulefill

\subsection*{Problem 7}
\textbf{Find an integrating factor of the form $x^n y^m$ and solve the equation.}
\[ (3xy^3 - 6y^2) \, dx + (4x^2y^2 - 9xy) \, dy = 0 \]

\textbf{Step 1: Multiply the equation by $x^n y^m$.}
\begin{align*}
    M' &= 3x^{n+1}y^{m+3} - 6x^ny^{m+2} \\
    N' &= 4x^{n+2}y^{m+2} - 9x^{n+1}y^{m+1}
\end{align*}

\textbf{Step 2: Force exactness ($M'_y = N'_x$).}
\begin{align*}
    M'_y &= 3(m + 3)x^{n+1}y^{m+2} - 6(m + 2)x^ny^{m+1} \\
    N'_x &= 4(n + 2)x^{n+1}y^{m+2} - 9(n + 1)x^ny^{m+1}
\end{align*}
Equate the coefficients of corresponding powers of $x$ and $y$:
1) For $x^{n+1}y^{m+2}$: $3(m + 3) = 4(n + 2) \implies 3m + 9 = 4n + 8 \implies 3m - 4n = -1$ \\
2) For $x^ny^{m+1}$: $-6(m + 2) = -9(n + 1) \implies 2(m + 2) = 3(n + 1) \implies 2m - 3n = -1$

\textbf{Step 3: Solve the system of equations for $m$ and $n$.}
\begin{align*}
    3m - 4n &= -1 \\
    2m - 3n &= -1 \implies 2m = 3n - 1 \implies m = \frac{3n - 1}{2}
\end{align*}
Substitute $m$ into the first equation:
\[ 3\left(\frac{3n - 1}{2}\right) - 4n = -1 \implies 9n - 3 - 8n = -2 \implies n = 1 \]
If $n = 1$, then $m = \frac{3(1)-1}{2} = 1$. The integrating factor is $\mu(x, y) = xy$.

\textbf{Step 4: Solve the newly exact equation.} Multiply original ODE by $xy$:
\[ (3x^2y^4 - 6xy^3) \, dx + (4x^3y^3 - 9x^2y^2) \, dy = 0 \]
Integrate $M'$ with respect to $x$:
\[ F(x, y) = \int (3x^2y^4 - 6xy^3) \, dx = x^3y^4 - 3x^2y^3 + g(y) \]
Differentiate with respect to $y$:
\[ \frac{\partial F}{\partial y} = 4x^3y^3 - 9x^2y^2 + g'(y) = 4x^3y^3 - 9x^2y^2 \implies g'(y) = 0 \implies g(y) = 0 \]
The solution is $x^3y^4 - 3x^2y^3 = C$.

\textbf{Step 5: Check for lost or gained solutions.} We multiplied by $xy$, which does not restrict the domain (it doesn't cause division by zero). Therefore, we didn't lose any solutions. (Checking $x = 0$ or $y = 0$ in the original shows they were solutions, and they remain solutions as $0 = 0$ still holds).

\textbf{Final Answers:}
\begin{itemize}
    \item Implicit solution: $x^3y^4 - 3x^2y^3 = C$
    \item Lost/Gained: \textbf{no solutions were lost}
\end{itemize}

\newpage
\section*{Part 3: In-Depth Analysis of Problems and Techniques}

\subsection*{A) Problem Types and General Approach}
The problems in this assignment fall into four distinct categories:

\textbf{1. Verification / Setup Only (Problems 1, 2):} These problems ask you to calculate $M_y$ and $N_x$ to classify the differential equation. The strategy is straightforward partial differentiation. If $M_y = N_x$, it is exact. If not, use the $\mu(x)$ or $\mu(y)$ formulas to see if an integrating factor exists.

\textbf{2. Exact Equations (Problem 6):} The most fundamental type. Once verified ($M_y = N_x$), the strategy is to integrate $M$ with respect to $x$ (adding $g(y)$), differentiate that result with respect to $y$, set it equal to $N$, and solve for $g'(y)$.

\textbf{3. Non-Exact Requiring 1D Integrating Factor (Problems 3, 4, 5):} Here, $M_y \neq N_x$. The approach requires calculating either $\frac{M_y - N_x}{N}$ or $\frac{N_x - M_y}{M}$. Whichever yields a single-variable expression dictates whether to use $\mu(x)$ or $\mu(y)$. Once $\mu$ is found and multiplied through, the equation falls back into the ``Exact Equation'' category.

\textbf{4. Non-Exact Requiring Algebraic Integrating Factor (Problem 7):} The integrating factor is explicitly given as a form $\mu = x^n y^m$. The strategy here relies heavily on algebra: multiply the literal terms $x^n y^m$ through the equation, compute the new mixed partials $M'_y$ and $N'_x$, and set the coefficients of corresponding polynomials equal to each other to form a system of linear equations.

\subsection*{B) Key Algebraic and Calculus Manipulations}
\begin{itemize}
    \item \textbf{Factoring to Isolate the Variable (Problems 1, 4):} When testing for an integrating factor $\mu(x)$, the expression $\frac{M_y - N_x}{N}$ initially looks messy (e.g., $\frac{x^{-1} - y^3}{xy^3 - 1}$). The critical algebraic trick is factoring out a term (in this case, $-x^{-1}$) from the numerator to perfectly cancel the denominator. Without this factorization skill, it is impossible to see that the term is a function of $x$ alone.
    
    \item \textbf{Exponential/Logarithmic Cancellation (Problems 1, 3, 4, 5):} Generating the integrating factor always involves computing $e^{\int p(x)dx}$. In all these problems, the integral yields a logarithm, such as $\int -\frac{3}{x} \, dx = -3\ln|x|$. The key manipulation is moving the coefficient into the exponent before canceling: $e^{-3\ln|x|} = e^{\ln(|x|^{-3})} = x^{-3}$.
    
    \item \textbf{Systems of Linear Equations from Polynomials (Problem 7):} By forcing $M'_y = N'_x$, we obtain an equation of polynomials. The profound algebraic concept used here is that two polynomials are identically equal if and only if their corresponding coefficients are equal. This converts a calculus problem into a linear algebra problem ($3m - 4n = -1$ and $2m - 3n = -1$).
    
    \item \textbf{Determining Gained/Lost Solutions:} Multiplying an ODE by a function like $x^{-3}$ implicitly assumes $x \neq 0$. If $x = 0$ was a valid solution originally, we ``lost'' it because our final closed-form equation will be undefined at $x = 0$. Conversely, multiplying by $x^2$ (as in Problem 5) can heal a singularity, creating a root at $x = 0$ that was undefined in the original equation, thus ``gaining'' a solution.
\end{itemize}

\newpage
\section*{Part 4: ``Cheatsheet'' and Tips for Success}

\subsection*{Summary of Most Important Formulas}
\begin{itemize}
    \item \textbf{Standard Form:} $M(x, y) \, dx + N(x, y) \, dy = 0$
    \item \textbf{Exactness Condition:} $M_y = N_x$
    \item \textbf{Implicit Solution Construction:} $F(x, y) = \int M \, dx + g(y) = C$
    \item \textbf{Finding Integrating Factors:}
    \begin{itemize}
        \item If $\frac{M_y - N_x}{N} = h(x)$, then $\mu(x) = e^{\int h(x) \, dx}$
        \item If $\frac{N_x - M_y}{M} = k(y)$, then $\mu(y) = e^{\int k(y) \, dy}$
    \end{itemize}
\end{itemize}

\subsection*{Tricks, ``Cheats'', and Shortcuts}
\begin{itemize}
    \item \textbf{The Denominator Clue:} When trying to figure out which integrating factor formula to use, look at $M$ and $N$. If $N$ is simpler (fewer terms), try $\frac{M_y - N_x}{N}$ first. If $M$ is simpler, try $\frac{N_x - M_y}{M}$. You want the simpler term in the denominator to make cancellation easier.
    \item \textbf{Integration Shortcut for Exact Equations:} To find $F(x, y)$, integrate $M$ with respect to $x$, and integrate $N$ with respect to $y$. Write down all the resulting terms, but \textit{do not write overlapping terms twice}. E.g., if $\int M \, dx = x^2y + x$ and $\int N \, dy = x^2y + y^2$, the final function is just the union of unique terms: $F(x, y) = x^2y + x + y^2 = C$.
    \item \textbf{Negative Sign Trick:} If your fraction for $\frac{M_y - N_x}{N}$ looks exactly like the denominator but with the terms swapped (e.g., $\frac{1 - y}{y - 1}$), factor out a $-1$ immediately to reveal the answer.
\end{itemize}

\subsection*{Common Pitfalls and Traps}
\begin{itemize}
    \item \textbf{Wrong Partial Derivatives:} The most common mistake is taking $M_x$ and $N_y$ instead of the cross-partials $M_y$ and $N_x$. Remember, $M$ is paired with $dx$, so you differentiate it with respect to $y$.
    \item \textbf{Forgetting the Constant $C$:} An exact equation yields an implicit solution $F(x, y) = C$. Writing $F(x, y) = 0$ or just leaving it as an expression without the ``$= C$'' will cost you full marks.
    \item \textbf{Botching the Integrating Factor Integral:} Remembering $e^{\int h(x) \, dx}$ but forgetting to use logarithmic properties correctly. E.g., $e^{-\ln x} \neq -x$. It is $x^{-1}$.
\end{itemize}

\newpage
\section*{Part 5: Conceptual Synthesis and The ``Big Picture''}

\subsection*{A) Thematic Connections}
The core theme of this topic is \textbf{Path Independence and Potential Functions}. If you have studied Multivariable Calculus, the equation $M \, dx + N \, dy = 0$ is functionally identical to finding the potential function of a conservative vector field $\vec{F} = \langle M, N \rangle$. When a vector field is conservative, the line integral $\int \vec{F} \cdot d\vec{r}$ is path-independent, and $\vec{F} = \nabla f$. Setting $M_y = N_x$ is exactly the 2D curl test: $\frac{\partial N}{\partial x} - \frac{\partial M}{\partial y} = 0$. Finding the implicit solution $F(x, y)$ is mathematically identical to finding the potential function $f(x, y)$ for the vector field!

\subsection*{B) Forward and Backward Links}
\textbf{Backward Links:} This builds heavily on the \textit{Chain Rule} for multivariable functions. The total differential $dF = F_x \, dx + F_y \, dy$ is the foundation of this entire chapter. If you understand how changes in $x$ and $y$ simultaneously affect $F$, exact equations make perfect logical sense.

\textbf{Forward Links:} Mastery of exact equations and integrating factors is vital for:
\begin{itemize}
    \item \textbf{Green's Theorem \& Line Integrals:} In Vector Calculus, evaluating line integrals over closed curves relies heavily on identifying whether the integrand forms an exact differential.
    \item \textbf{Thermodynamics:} In physics, state variables (like Entropy or Internal Energy) are ``exact differentials'' (e.g., $dU = T \, dS - P \, dV$), whereas heat and work are path-dependent (inexact differentials). Integrating factors (like $1/T$) are literally used to turn inexact heat differentials into exact entropy differentials!
\end{itemize}

\section*{Part 6: Real-World Application and Modeling}

\subsection*{A) Concrete Scenarios in Financial Modeling}
\begin{enumerate}
    \item \textbf{Interest Rate Modeling (The Vasicek Model):} In quantitative finance, the evolution of interest rates is often modeled using a mean-reverting stochastic differential equation. To find the expected value of the interest rate at a future time $t$, quants must solve a deterministic linear first-order ODE. This equation invariably requires an integrating factor to solve, ensuring the model accurately predicts how rates are ``pulled'' back to a long-term average.
    \item \textbf{Continuous-Time Wealth Equations:} When analyzing a portfolio that continuously pays dividends or consumes cash flow, the wealth equation takes the form of a first-order linear ODE. An integrating factor is used to discount future cash flows back to present value, transforming the differential equation into an easily solvable exact equation.
\end{enumerate}

\subsection*{B) Model Problem Setup: Expected Interest Rates}
Let's model the \textbf{Vasicek Interest Rate} expected value. Suppose the short-term interest rate $r(t)$ follows a mean-reverting process. By taking the expected value, the stochastic noise cancels out, and we are left with a deterministic model for the expected interest rate $y(t) = \mathbb{E}[r(t)]$:
\[ \frac{dy}{dt} = a(b - y) \]
where $a$ is the speed of mean reversion, and $b$ is the long-term mean level.

\textbf{Setup:} We can rewrite this as:
\[ dy = a(b - y) \, dt \implies dy + a(y - b) \, dt = 0 \implies (ay - ab) \, dt + (1) \, dy = 0 \]
Here, the variables are $t$ and $y$.
\begin{align*}
    M(t, y) &= ay - ab \implies M_y = a \\
    N(t, y) &= 1 \implies N_t = 0
\end{align*}
Since $M_y \neq N_t$, the equation is not exact. To find the expected interest rate over time, a quant would compute the integrating factor:
\[ \frac{M_y - N_t}{N} = \frac{a - 0}{1} = a \]
This gives an integrating factor of $\mu(t) = e^{\int a \, dt} = e^{at}$. Multiplying the ODE by this factor allows the quant to find the explicit formula for $y(t)$, showing how rates exponentially decay toward $b$.

\newpage
\section*{Part 7: Common Variations and Untested Concepts}

The homework primarily focused on Exact Equations and arbitrary integrating factors. However, the standard curriculum for this topic usually explicitly includes two standard forms that act as ``shortcuts.''

\subsection*{1. The Standard First-Order Linear ODE}
Often, a problem is presented in the form:
\[ \frac{dy}{dx} + P(x)y = Q(x) \]
This is just a specific rearrangement of $M \, dx + N \, dy = 0$. For this exact form, the integrating factor is \textit{always} $\mu(x) = e^{\int P(x) \, dx}$. Multiplying the whole equation by $\mu(x)$ guarantees that the left side becomes the exact derivative of a product: $\frac{d}{dx}[\mu(x)y]$.

\textbf{Example Problem:} Solve $\frac{dy}{dx} + \frac{2}{x}y = x^2$ for $x > 0$.
\textbf{Solution:} Here $P(x) = \frac{2}{x}$.
\[ \mu(x) = e^{\int \frac{2}{x} \, dx} = e^{2\ln|x|} = x^2 \]
Multiply the entire equation by $x^2$:
\[ x^2\frac{dy}{dx} + 2xy = x^4 \]
The left side collapses by the reverse product rule:
\[ \frac{d}{dx}[x^2y] = x^4 \]
Integrate both sides with respect to $x$:
\[ x^2y = \frac{x^5}{5} + C \implies y = \frac{x^3}{5} + Cx^{-2} \]

\subsection*{2. Bernoulli Differential Equations}
A Bernoulli equation is a non-linear ODE of the form:
\[ \frac{dy}{dx} + P(x)y = Q(x)y^n \]
It can be transformed into a linear ODE (and thus solved with an integrating factor) using the substitution $v = y^{1-n}$.

\textbf{Example Problem:} Solve $x\frac{dy}{dx} + y = x^2y^2$. 
\textbf{Solution:} Divide by $x$: $\frac{dy}{dx} + \frac{1}{x}y = xy^2$. This is Bernoulli with $n = 2$. Let $v = y^{1-2} = y^{-1}$. Then $\frac{dv}{dx} = -y^{-2}\frac{dy}{dx}$. Divide the standard equation by $y^2$:
\[ y^{-2}\frac{dy}{dx} + \frac{1}{x}y^{-1} = x \]
Substitute $v$ and $dv/dx$:
\[ -\frac{dv}{dx} + \frac{1}{x}v = x \implies \frac{dv}{dx} - \frac{1}{x}v = -x \]
Now it is a standard linear ODE! The integrating factor is $\mu(x) = e^{\int -1/x \, dx} = x^{-1}$. Multiply by $x^{-1}$:
\[ x^{-1}\frac{dv}{dx} - x^{-2}v = -1 \implies \frac{d}{dx}[x^{-1}v] = -1 \]
Integrate: $x^{-1}v = -x + C \implies v = -x^2 + Cx$. Substitute back $y$: $y^{-1} = -x^2 + Cx \implies y = \frac{1}{Cx - x^2}$.

\newpage
\section*{Part 8: Advanced Diagnostic Testing: ``Find the Flaw''}

Below are 5 flawed solutions to ODE problems. Act as the grader: find the error, explain why it is wrong in one sentence, and provide the correct step and final answer.

\subsection*{Flawed Problem 1}
\textbf{Problem:} Determine if $(2x + y) \, dx + (x - 2y) \, dy = 0$ is exact.
\\
\textbf{Flawed Solution:} Let $M = 2x + y$ and $N = x - 2y$. Take the partial derivative of $M$ with respect to $x$: $M_x = 2$. Take the partial derivative of $N$ with respect to $y$: $N_y = -2$. Since $M_x \neq N_y$, the equation is not exact.

\textbf{Grader Correction:}
\begin{itemize}
    \item \textbf{The Flaw:} The student took the wrong partial derivatives; they computed $M_x$ and $N_y$ instead of the required cross-partial derivatives $M_y$ and $N_x$.
    \item \textbf{Correct Step:} Take the partial derivative of $M$ with respect to $y$ ($M_y = 1$) and $N$ with respect to $x$ ($N_x = 1$). 
    \item \textbf{Final Answer:} Since $M_y = N_x = 1$, the equation \textbf{is exact}.
\end{itemize}

\subsection*{Flawed Problem 2}
\textbf{Problem:} Solve the exact equation $2xy \, dx + x^2 \, dy = 0$.
\\
\textbf{Flawed Solution:} $M_y = 2x$ and $N_x = 2x$. Exact. Integrate $M$ with respect to $x$: $\int 2xy \, dx = x^2y$. Integrate $N$ with respect to $y$: $\int x^2 \, dy = x^2y$. Add the two results together to get $F(x, y)$: $F(x, y) = x^2y + x^2y = 2x^2y$. Solution is $2x^2y = C$.

\textbf{Grader Correction:}
\begin{itemize}
    \item \textbf{The Flaw:} The student incorrectly added the integral results rather than taking the union of the unique terms to form the potential function $F(x, y)$.
    \item \textbf{Correct Step:} When combining the partial integrals $\int M \, dx = x^2y$ and $\int N \, dy = x^2y$, overlapping terms should only be counted once. Thus, $F(x,y) = x^2y$.
    \item \textbf{Final Answer:} The correct implicit solution is $x^2y = C$.
\end{itemize}

\subsection*{Flawed Problem 3}
\textbf{Problem:} Find the integrating factor for $(3xy + y^2) \, dx + (x^2 + xy) \, dy = 0$.
\\
\textbf{Flawed Solution:} $M_y = 3x + 2y$, $N_x = 2x + y$. Not exact. Check $\frac{M_y - N_x}{N}$:
\[ \frac{(3x + 2y) - (2x + y)}{x^2 + xy} = \frac{x + y}{x(x + y)} = \frac{1}{x} \]
This is a function of $x$. Thus, $\mu(y) = e^{\int 1/x \, dy} = e^{y/x}$.

\textbf{Grader Correction:}
\begin{itemize}
    \item \textbf{The Flaw:} The student correctly found that the fraction depends only on $x$, but erroneously set up the formula to find $\mu(y)$ and integrated with respect to $y$.
    \item \textbf{Correct Step:} Because the ratio $\frac{M_y - N_x}{N} = \frac{1}{x}$ is a function of $x$ alone, we must compute $\mu(x) = e^{\int \frac{1}{x} \, dx}$.
    \item \textbf{Final Answer:} $\mu(x) = e^{\ln|x|} = x$.
\end{itemize}

\subsection*{Flawed Problem 4}
\textbf{Problem:} Find the integrating factor $\mu(x)$ if $\frac{M_y - N_x}{N} = -\frac{2}{x}$.
\\
\textbf{Flawed Solution:} $\mu(x) = e^{\int -\frac{2}{x} \, dx} \implies \mu(x) = e^{-2\ln|x|}$. Using logarithm rules, $e^{-2\ln|x|} = -2x$. So $\mu(x) = -2x$.

\textbf{Grader Correction:}
\begin{itemize}
    \item \textbf{The Flaw:} The student fundamentally misused logarithm and exponential properties; a coefficient in front of a natural log moves to the exponent of the argument, it does not become a standard coefficient.
    \item \textbf{Correct Step:} Apply the power rule for logarithms first: $e^{-2\ln|x|} = e^{\ln(|x|^{-2})}$. 
    \item \textbf{Final Answer:} $\mu(x) = x^{-2}$ or $\frac{1}{x^2}$.
\end{itemize}

\subsection*{Flawed Problem 5}
\textbf{Problem:} Solve $x\frac{dy}{dx} + 2y = x^2$ using an integrating factor.
\\
\textbf{Flawed Solution:} The equation is of the form $\frac{dy}{dx} + P(x)y = Q(x)$ where $P(x) = 2$. The integrating factor is $\mu(x) = e^{\int 2 \, dx} = e^{2x}$. Multiply by $e^{2x}$: $xe^{2x}y' + 2e^{2x}y = x^2e^{2x}$. Integrate both sides...

\textbf{Grader Correction:}
\begin{itemize}
    \item \textbf{The Flaw:} The student identified $P(x)$ before putting the linear ODE into standard form (which requires the coefficient of $\frac{dy}{dx}$ to be exactly $1$).
    \item \textbf{Correct Step:} First divide the entire equation by $x$ to get $\frac{dy}{dx} + \frac{2}{x}y = x$. Then $P(x) = \frac{2}{x}$, which gives $\mu(x) = e^{\int \frac{2}{x} \, dx} = e^{2\ln|x|} = x^2$.
    \item \textbf{Final Answer:} Multiply standard form by $x^2$ to get $x^2y' + 2xy = x^3 \implies \frac{d}{dx}[x^2y] = x^3 \implies x^2y = \frac{x^4}{4} + C \implies y = \frac{x^2}{4} + Cx^{-2}$.
\end{itemize}

\end{document}